\documentclass[12pt,a4paper]{article}
\usepackage[utf8]{inputenc}
\usepackage[T1]{fontenc}
\usepackage[russian]{babel}
\usepackage{amsmath}
\usepackage{amssymb}
\usepackage{graphicx}
\usepackage{tikz}
\usepackage{geometry}
\geometry{a4paper, margin=2.5cm}

\usetikzlibrary{arrows, decorations.markings, calc}

\title{Классическая модель спина электрона и альтернативная модель прецессирующего магнитного диполя}
\author{Алексей Дроздов}
\date{\today}

\begin{document}

\maketitle

\begin{abstract}
В работе рассматривается классическая проблема объяснения магнитного момента электрона через вращение заряженного шара и предлагается альтернативная модель, основанная на прецессии магнитного диполя. Показано, что классическая модель требует сверхсветовых скоростей вращения, тогда как модель прецессирующего магнитного диполя дает физически осмысленные результаты.
\end{abstract}

\section{Классическая модель: вращающийся заряженный шар}

\subsection{Постановка задачи}

Рассмотрим классическую модель электрона как равномерно заряженную сферу радиуса $R$, вращающуюся с угловой скоростью $\omega$. Цель -- объяснить наблюдаемый магнитный момент электрона (магнетон Бора $\mu_B$) через механическое вращение заряда.

\begin{figure}[h]
\centering
\begin{tikzpicture}[scale=1.2]
% Сфера
\draw[thick] (0,0) circle (2cm);
\draw[dashed] (0,0) ellipse (2cm and 0.5cm);
\draw[dashed] (0,0) ellipse (0.5cm and 2cm);

% Ось вращения
\draw[thick,->] (0,-2.5) -- (0,2.5);
\node at (0.3,2.7) {$\omega$};

% Заряд на поверхности
\node at (-1.8,-1.6) {$e^-$};

% Одна стрелка вращения на экваторе
\draw[thick,->,blue] (1.5,0) arc (0:120:1.5cm and 0.4cm);


% Радиус
\draw[thick,->] (0,0) -- (1.4,1.4);
\node at (0.9,0.9) {$R$};
\end{tikzpicture}
\caption{Классическая модель: вращающаяся заряженная сфера}
\label{fig:classical}
\end{figure}


\subsection{Магнитный момент вращающейся сферы}

Поверхностная плотность заряда:
\begin{equation}
\sigma = \frac{e}{4\pi R^2}
\end{equation}

Элемент заряда на широте $\theta$:
\begin{equation}
dq = \sigma \cdot 2\pi R \sin\theta \cdot R d\theta = \frac{e}{2} \sin\theta \, d\theta
\end{equation}

Этот элемент создает ток:
\begin{equation}
dI = \frac{dq \cdot \omega}{2\pi} = \frac{e\omega}{4\pi} \sin\theta \, d\theta
\end{equation}

Магнитный момент элемента (площадь кольца $\pi R^2 \sin^2\theta$):
\begin{equation}
dM = \frac{dI \cdot \pi R^2 \sin^2\theta}{c} = \frac{e\omega R^2}{4c} \sin^3\theta \, d\theta
\end{equation}

Интегрируем по всей сфере:
\begin{equation}
M = \frac{e\omega R^2}{4c} \int_0^\pi \sin^3\theta \, d\theta
\end{equation}

Вычисляем интеграл:
\begin{equation}
\int_0^\pi \sin^3\theta \, d\theta = \int_0^\pi \sin\theta(1-\cos^2\theta) \, d\theta = \left[-\cos\theta + \frac{\cos^3\theta}{3}\right]_0^\pi = \frac{4}{3}
\end{equation}

Получаем магнитный момент:
\begin{equation}
\boxed{M = \frac{e\omega R^2}{3c}}
\label{eq:classical_moment}
\end{equation}

\subsection{Классический радиус электрона}

Из условия равенства энергии электростатического поля энергии покоя:
\begin{equation}
W = \frac{e^2}{2R} = m_e c^2
\end{equation}

Отсюда классический радиус:
\begin{equation}
R = \frac{e^2}{2m_e c^2} = \frac{(4.8 \times 10^{-10})^2}{2 \times 9.1 \times 10^{-28} \times (3 \times 10^{10})^2} \approx 1.4 \times 10^{-13} \text{ см}
\end{equation}

\subsection{Проблема сверхсветовой скорости}

Приравниваем классический магнитный момент к магнетону Бора:
\begin{equation}
M_B = \frac{e\omega R^2}{3c} = 9.274 \times 10^{-21} \text{ эрг/Гс}
\end{equation}

Выражаем линейную скорость на экваторе $v = \omega R$:
\begin{equation}
v = \omega R = \frac{3c M_B}{e R}
\end{equation}

Подставляем численные значения:
\begin{equation}
v = \frac{3 \times 3 \times 10^{10} \times 9.274 \times 10^{-21}}{4.8 \times 10^{-10} \times 1.4 \times 10^{-13}}
\end{equation}

\begin{equation}
v = \frac{8.35 \times 10^{-10}}{6.72 \times 10^{-23}} \approx 1.24 \times 10^{13} \text{ см/с}
\end{equation}

Отношение к скорости света:
\begin{equation}
\boxed{\frac{v}{c} \approx 414}
\end{equation}

\textbf{Вывод:} Классическая модель требует, чтобы поверхность электрона вращалась со скоростью, превышающей скорость света в 414 раз, что физически невозможно.

\section{Альтернативная модель: прецессирующий магнитный диполь}

\subsection{Постановка задачи}

Рассмотрим электрон как систему двух магнитных зарядов $+q_m$ и $-q_m$, разделенных расстоянием $2R$ и прецессирующих вокруг общей оси.

\begin{figure}[h]
\centering
\begin{tikzpicture}[scale=1.2]
% Верхний магнитный заряд
\draw[thick] (0,1.5) ellipse (1.5cm and 0.4cm);
\draw[thick,->,blue] (1.5,1.5) arc (0:180:1.5cm and 0.4cm);
\draw[thick,->,blue] (-1.5,1.5) arc (180:360:1.5cm and 0.4cm);
\filldraw[red] (1.5,1.5) circle (0.15);
\node at (2,1.7) {$+q_m$};

% Нижний магнитный заряд
\draw[thick] (0,-1.5) ellipse (1.5cm and 0.4cm);
\draw[thick,->,blue] (1.5,-1.5) arc (0:180:1.5cm and 0.4cm);
\draw[thick,->,blue] (-1.5,-1.5) arc (180:360:1.5cm and 0.4cm);
\filldraw[blue] (-1.5,-1.5) circle (0.15);
\node at (-2,-1.7) {$-q_m$};

% Ось прецессии
\draw[thick,->] (0,-2.5) -- (0,2.5);
\node at (0.3,2.7) {$\omega$};

% Расстояние
\draw[thick,<->] (2,1.5) -- (2,-1.5);
\node at (2.8,0) {$2R \cdot cos \varphi$};

% Диагональная линия связи 2R
\draw[thick,<->] (-1.5,-1.5) -- (1.5,1.5);
\node at (1.3,0.5) {$2R$};

% Центральная точка (центр диполя)
\filldraw[black] (0, 0) circle (0.08);

% Стрелки прецессии
\draw[thick,->,green!60!black] (0.5,1.5) arc (20:160:0.5cm);
\draw[thick,->,green!60!black] (-0.5,-1.5) arc (200:340:0.5cm);
\end{tikzpicture}

\caption{Модель прецессирующего магнитного диполя: два магнитных заряда $+q_m$ и $-q_m$ прецессируют по кольцам радиуса $R \cdot sin \varphi$, расстояние между зарядами $2R$}
\label{fig:precession}
\end{figure}

\subsection{Магнитный момент диполя}

Магнитный момент системы двух магнитных зарядов:
\begin{equation}
M_B = 2R \cdot q_m
\label{eq:magnetic_dipole}
\end{equation}

Из условия равенства магнетону Бора находим величину магнитного заряда:
\begin{equation}
q_m = \frac{M_B}{2R} = \frac{9.274 \times 10^{-21}}{2 \times 1.4 \times 10^{-13}} = 3.312 \times 10^{-8} \text{ СГСМ}
\end{equation}

\subsection{Связь движения магнитного заряда с потоком электрического поля}

\subsubsection{Задача ``наизнанку''}

В классической электродинамике вращение электрического заряда создает магнитное поле. Рассмотрим обратную задачу: вращение магнитного заряда создает электрическое поле.

Из модифицированных уравнений Максвелла с магнитными зарядами:
\begin{equation}
\oint \vec{E} \cdot d\vec{l} = -\frac{4\pi}{c} I_m
\end{equation}

где $I_m$ -- магнитный ток.

\subsubsection{Магнитный ток}

Магнитный заряд $q_m$, вращающийся по кольцу радиуса $R$ с угловой скоростью $\omega$:
\begin{equation}
I_m = \frac{q_m \omega}{2\pi}
\end{equation}

\subsubsection{Электрическое поле}

Для кольца радиуса $R$:
\begin{equation}
E \cdot 2\pi R = \frac{4\pi}{c} \cdot \frac{q_m \omega}{2\pi}
\end{equation}

\begin{equation}
E = \frac{q_m \omega}{\pi c R}
\label{eq:e_field}
\end{equation}

\subsubsection{Поток электрического поля через кольцо}

\begin{equation}
\Phi_E^{ring} = E \cdot \pi R^2 = \frac{q_m \omega R}{c}
\end{equation}

\subsubsection{Учет двух магнитных зарядов}

Оба магнитных заряда ($+q_m$ и $-q_m$) вращаются в одном направлении, создавая электрические поля, направленные одинаково (знак заряда компенсируется направлением вращения):

\begin{equation}
\Phi_E^{total} = 2 \cdot \frac{q_m \omega R}{c} = \frac{2q_m \omega R}{c}
\label{eq:total_flux}
\end{equation}

\subsection{Сравнение с потоком от электрического заряда}

Электрический заряд $e$ на сфере радиуса $R$ создает поток:
\begin{equation}
\Phi_E^{sphere} = 4\pi e
\end{equation}

Приравниваем потоки для получения эффективного электрического заряда:
\begin{equation}
4\pi e = \frac{2q_m \omega R}{c}
\end{equation}

Отсюда получаем основное уравнение:
\begin{equation}
\boxed{\omega R = \frac{2\pi c e}{q_m}}
\label{eq:main}
\end{equation}

Или в безразмерном виде:
\begin{equation}
\frac{\omega R}{c} = \frac{2\pi e}{q_m}
\end{equation}

\subsection{Численный расчет}

Подставляем значения:
\begin{align}
\omega R &= \frac{2\pi \times 3 \times 10^{10} \times 4.8 \times 10^{-10}}{3.312 \times 10^{-8}} \\
&= \frac{2\pi \times 14.4}{3.312 \times 10^{-8}} \\
&\approx 2.73 \times 10^9 \text{ см/с}
\end{align}

Отношение к скорости света:
\begin{equation}
\boxed{\frac{\omega R}{c} = \frac{2\pi e}{q_m} = \frac{2\pi \times 4.8 \times 10^{-10}}{3.312 \times 10^{-8}} \approx 0.091}
\end{equation}

\subsection{Альтернативный вывод через магнитный ток}

Можно также рассмотреть задачу через магнитный ток. Вращающийся магнитный заряд создает магнитный ток:

\begin{equation}
J_m = \frac{q_m \omega_m}{2\pi}
\end{equation}

Этот магнитный ток через уравнения Максвелла создает электрическое поле, эквивалентное полю от электрического тока:
\begin{equation}
J_m = c \frac{e}{\pi R}
\end{equation}

Отсюда:
\begin{equation}
\frac{q_m \omega_m}{2\pi} = c \frac{e}{\pi R}
\end{equation}

\begin{equation}
\omega_m R = 2c \frac{e}{q_m}
\end{equation}

Численно:
\begin{equation}
\frac{\omega R}{c} = \frac{2e}{q_m} = \frac{2 \times 4.8 \times 10^{-10}}{3.312 \times 10^{-8}} \approx 2.89 \times 10^{-2}
\end{equation}

\section{Сравнение моделей}

\begin{table}[h]
\centering
\begin{tabular}{|l|c|c|}
\hline
\textbf{Параметр} & \textbf{Классическая модель} & \textbf{Модель прецессии} \\
\hline
Исходный объект & Вращающийся заряд $e$ & Прецессирующий диполь $q_m$ \\
\hline
Формула момента & $M = \frac{e\omega R^2}{3c}$ & $M = 2R q_m$ \\
\hline
Требуемая скорость & $v \approx 414c$ & $v \approx 0.03c - 0.09c$ \\
\hline
Физическая осмысленность & Нарушает СТО & Согласуется с СТО \\
\hline
\end{tabular}
\caption{Сравнение классической и альтернативной моделей}
\label{tab:comparison}
\end{table}

\section{Выводы}

\begin{enumerate}
\item Классическая модель вращающегося заряженного шара требует сверхсветовых скоростей ($v \approx 414c$) для объяснения магнитного момента электрона, что делает её физически несостоятельной.

\item Модель прецессирующего магнитного диполя предлагает альтернативное объяснение: электрон представляет собой систему двух магнитных зарядов, прецессирующих вокруг общей оси.

\item Вращение магнитных зарядов создает эффективный электрический заряд через поток электрического поля, что является задачей ``наизнанку'' по отношению к классической задаче о магнитном потоке токового кольца.

\item Требуемая скорость прецессии составляет всего $\sim 3-9\%$ от скорости света, что физически осмысленно и не противоречит специальной теории относительности.

\item Полученный результат подтверждает гипотезу о том, что ``магнитная природа'' электрона является более фундаментальной, чем ``электрическая природа''.
\end{enumerate}

\begin{thebibliography}{9}
\bibitem{classical} Абрагам М. Теория электричества. -- М.: ОНТИ, 1934.
\bibitem{dirac} Дирак П. Квантованная сингулярность в электромагнитном поле. // Труды Королевского общества. -- 1931. -- Т. A133. -- С. 60-72.
\bibitem{magnetic_charge} Джексон Дж. Классическая электродинамика. -- М.: Мир, 1965.
\end{thebibliography}

\end{document}
